% Hoja 2
\setcounter{ejercicio}{0}

% 1
\ejercicio{Calcula el subespacio proyectivo de $\mathbb{P}^3_{\mathbb{R}}$ generado por los puntos:\\
$\underbrace{(0 : 1 : 0 : 0)}_{\text{A}}, \underbrace{(2 : -1 : 0 : 1)}_{\text{B}}, \underbrace{(-4 : 1 : 0 : 2)}_{\text{C}}$}
{
    Primero veamos la dimension de $V = (A, B, C)$:
    \[
        \dim(V(A,B,C)) =
        rg \begin{pmatrix}
            0 & 1 & 0 & 0 \\
            2 & -1 & 0 & 1 \\
            -4 & 1 & 0 & 2
        \end{pmatrix} - 1 = 2
    \]
    En general si $X \subset \mathbb{P}^n$ y $\dim(X) = k$, entonces $X$ tiene $n-k$ ecuaciones implícitas. Así obtengamos las ecuaciones implícitas de $X$:
    \[
        \begin{pmatrix}
            x_0 & x_1 & x_2 & x_3 \\
            0 & 1 & 0 & 0 \\
            2 & -1 & 0 & 1 \\
            -4 & 1 & 0 & 2
        \end{pmatrix} = 0 \implies V(A,B,C) : X_2 = 0
    \]
}

% 2
\ejercicio{En $\mathbb{P}^2_{\mathbb{C}}$ halla:
\begin{itemize}
    \item[a)] La ecuación de la recta que pasa por los puntos $(i : -1 : -i)$ y $(1 : 1 + i : 2)$
    \item[b)] El punto de intersección de las rectas $r : ix_0 + x_1 - ix_2 = 0$ y $s : x_0 + (1+i)x_1 + 2x_2 = 0$
\end{itemize}
}{
    \begin{itemize}
        \item [a)] La ecuación implicita de la recta que pasa por los puntos $A = (i : -1 : -i)$ y $B = (1 : 1 + i : 2)$ es:
        \[
            \begin{vmatrix}
                x_0 & x_1 & x_2 \\
                i & -1 & -i \\
                1 & 1 + i & 2
            \end{vmatrix} = 0 \implies (- 3 + i)x_0 + -3ix_1 + ix_2 = 0
        \]
        \item [b)] El punto de intersección de las rectas $r$ y $s$ es:
        \[
        \left\{
            \begin{array}{l}
                ix_0 + x_1 - ix_2 = 0 \\
                x_0 + (1+i)x_1 + 2x_2 = 0
            \end{array}
        \right.
        \iff
        \left\{
            \begin{array}{l}
                x_1 = i(x_2 - x_0) \\
                x_0 + (1+i)(i(x_2 - x_0)) + 2x_2 = 0
            \end{array}
        \right.
        \iff
        \]
        \[
        \left\{
            \begin{array}{l}
                x_1 = i(x_2 - x_0) \\
                x_0 + (i - 1)(x_2 - x_0) + 2x_2 = 0
            \end{array}
        \right.
        \iff
        (2-i)x_0 + (i+1)x_2 = 0 \implies x_2 = \frac{-(2-i)}{(1+i)}x_0
        \]
    \end{itemize}
}

% 3
\ejercicio{}{}

% 4
\ejercicio{}{}

% 5
\ejercicio{}{}

% 6
\ejercicio{}{}

% 7
\ejercicio{}{}

% 8
\ejercicio{}{}

% 9
\ejercicio{}{}

% 10
\ejercicio{}{}